\documentclass[10pt,a4paper]{amsart}
\usepackage[margin=19mm]{geometry}
\usepackage{amsmath,amssymb,mathtools}
\usepackage[scr=rsfs,cal=euler]{mathalfa}
\usepackage{xfrac}
\usepackage[unicode,hidelinks,bookmarks=false]{hyperref}
\newcommand{\ie}{i.e.\ }
\newcommand{\cf}{cf.\ }
\newcommand{\resp}{respectively\ }
\newcommand{\Subsection}{\subsection}
\providecommand{\nobreakdash}{\nobreak\hbox{-}}
\setlength{\emergencystretch}{2em}
\multlinegap0pt
\usepackage{bm}
\usepackage{bbm}
\usepackage{dsfont}
\usepackage{xspace}
\usepackage{cancel}
\usepackage{mathtools}
\usepackage{amsthm}
\makeatletter
\@ifundefined{thm}{\newtheorem{thm}{Thm}}{}
\@ifundefined{prop}{\newtheorem{prop}[thm]{Proposition}}{}
\makeatother
\title[Equidistant versus bipartite ground states for 1D classical ]{Equidistant versus bipartite ground states for 1D classical fluids at fixed particle density}
\author{Laurent B\'etermin, Ladislav \v{S}amaj, Igor Trav\v{e}nec}
\date{Source: arXiv:2502.16639v2 (CC BY 4.0)}
\begin{document}
\maketitle

\noindent\emph{Source and licence.} Equidistant versus bipartite ground states for 1D classical fluids at fixed particle density. Version arXiv:2502.16639v2 (CC BY 4.0), distributed under the Creative Commons Attribution 4.0 International licence, as recorded in the arXiv licence metadata for this exact version and shown on the abstract page https://arxiv.org/abs/2502.16639v2. Full rights notice: licenses/arxiv-2502.16639-CC-BY-4.0.txt.\par\smallskip

\noindent\textit{Editorial scope.} Counted unit: the Prop whose source label is \texttt{adel} in the article (source0.tex lines 1227--1243, source label \texttt{adel}), delivered together with the article's own complete original proof of it (source0.tex lines 1244--1292, one \texttt{proof} environment of 49 lines). No mathematics is written, repaired or re-derived: statement and proof are the article's own text line for line. Required context retained uncounted: none. Every definition, display and label of those lines is retained unchanged. Omitted from this package and not redistributed: 1--1226, 1293--1760. Nothing inside a retained excerpt is omitted or altered: every retained body is delivered exactly as the article's lines are written, so no omission receipt of any kind accompanies this package. Citations: the retained excerpts carry no citation command at all, so no bibliography entry of the article is transcribed and no reference is redistributed. Reference adaptation: none was needed and none was made; every reference inside the delivered unit resolves inside it, so no unresolved reference or citation is produced. Printed slips kept literal: no printed word, symbol, hypothesis or number of the counted statement or of its proof was altered. Layout and class: the article's own class is \texttt{iopart} with packages including graphics, graphicx, amssymb, amsthm, color; the wrapper is the lane's pinned \texttt{amsart} editorial wrapper at 10pt on A4, which restates the theorem-like environments the retained text needs and adds the article's own macro definitions that the retained text needs (none), with extra wrapper packages (bm, bbm, dsfont, xspace, cancel, mathtools). The delivered numbering is the unit's own. Assets: no retained span contains a figure environment, a graphics-inclusion command, a PDF-page inclusion, a table or a TikZ picture; no image file is copied into this package, the package declares no asset dependency and no image file is redistributed with it. Licence: version arXiv:2502.16639v2 of this article is distributed under the Creative Commons Attribution 4.0 International licence, as recorded in the arXiv licence metadata for this exact version and shown on the abstract page https://arxiv.org/abs/2502.16639v2. A full rights notice is delivered at licenses/arxiv-2502.16639-CC-BY-4.0.txt. Characters: every retained excerpt is pure ASCII, so the delivered engine, XeLaTeX with its default Latin Modern text font, reports no missing character.
\begin{prop}[\textbf{Bipartite phase and the asymptotics $A\to\infty$}]
The region of the bipartite phase is defined by $A>A_c^{nm}$:  
\begin{itemize}
\item[(1)]
for a given $A>A_c^{nm}$, there exists a nontrivial value of the parameter
$\Delta(A)\ne 1$ given by the equation
\begin{equation} \label{adel}
A^{m-n}=\frac{2^{n-m}[\zeta(m+1,1-\delta)-\zeta(m+1,\delta)]}{
\zeta(n+1,1-\delta)-\zeta(n+1,\delta)},\quad \delta=\frac{1}{1+\Delta}.
\end{equation}
\item[(2)] as $A\to +\infty$, we have
\begin{equation} \label{deltaas}
\Delta = 2A - 1 + o(1)
\end{equation}
regardless of the values of the LJ parameters $(n,m)$.
\end{itemize}
\end{prop}

\begin{proof}
The value of $\Delta$ is determined by the minimization condition for
the energy    
\begin{equation}
\frac{\partial}{\partial\Delta} E^{\rm bip}_{nm}(A,\Delta) = 0
\end{equation}
which gives
\begin{eqnarray}
\frac{1}{(2 A)^n} \left[ \zeta\left( n+1,\frac{1}{1+\Delta}\right) -
\zeta\left( n+1,\frac{\Delta}{1+\Delta}\right) \right] & & \nonumber \\
- \frac{1}{(2 A)^m} \left[ \zeta\left( m+1,\frac{1}{1+\Delta}\right) -
\zeta\left( m+1,\frac{\Delta}{1+\Delta}\right) \right] & = & 0 .   
\end{eqnarray}
It is clear that $\Delta=1$, which corresponds to the equidistant ground state,
always satisfies this equation; it provides the true energy minimum in the
region $0<A<A_c^{nm}$.
On the other hand, the trivial $\Delta=1$ corresponds to the energy maximum
in the bipartite region $A>A_c^{nm}$ where the nontrivial solution (\ref{adel})
with $\Delta>1$ provides the minimum of the energy.\\
We remark that in the limit $A\to +\infty$ it follows that, since $n>m$,
$$
\lim_{A\to +\infty} \frac{\zeta(m+1,1-\delta(A))-\zeta(m+1,\delta(A))}{
\zeta(n+1,1-\delta(A))-\zeta(n+1,\delta(A))}=0,\quad \delta(A)
=\frac{1}{1+\Delta(A)}
$$
which implies that $\displaystyle\lim_{A\to +\infty} \delta(A)=0$
(if not, the above limit cannot be $0$), and thus
$\displaystyle\lim_{A\to +\infty} \Delta(A)=+\infty$.
Therefore, in this $A\to +\infty$ regime, one can expand
the right-hand side of equation (\ref{adel}) in powers of small $\delta$
by using the expansion formulas
\begin{eqnarray} 
\zeta(m+1,\delta) = \frac{1}{\delta^{m+1}}+\zeta(m+1)-(m+1)\zeta(m+2)\delta
+ O\left( \delta^2\right) , \nonumber \\
\zeta(m+1,1-\delta) =  \zeta(m+1)+ (m+1) \zeta(m+1)\delta
+ O\left( \delta^2\right) , \label{zetaexp} 
\end{eqnarray}
implying
\begin{equation}
\zeta(m+1,1-\delta) - \zeta(m+1,\delta) = - \frac{1}{\delta^{m+1}}
+ 2 (m+1) \zeta(m+2) \delta + O\left( \delta^2\right) .
\end{equation}
The relation (\ref{adel}) thus yields
\begin{equation}
\left( \frac{2A}{1+\Delta} \right)^{m-n} = 1 +
O\left( \frac{1}{\Delta^{(3+m)}} \right)
\end{equation}  
which shows (\ref{deltaas}).
\end{proof}

\end{document}
